\begin{figure}[htbp]
\centering
\begin{tikzpicture}[font=\small,>=Stealth]

\draw[gray,thick] (-1,-.3) .. controls (-.6,0) and (-.3,0) .. (0,0)
 .. controls (1,0) and (2,0) .. (3.3,.7);
\fill (0,0) circle (2pt) node[below left] {$\gamma(t)$};
\draw[->,very thick,blue] (0,0)--(3,0) node[below] {$\gamma'(t)\Delta t$};
\draw[->,very thick,red!70!black] (0,0)--(2,2) node[above] {$F(\gamma(t))$};
\draw[dashed] (2,2)--(2,0);
\draw[->,thick,red!70!black] (0,.13)--(2,.13) node[midway,above] {$F_{\parallel}\tau$};
\draw (1.8,0)--(1.8,.2)--(2,.2);
\node[align=left,text width=6cm,anchor=west] at (4,1) {
$\tau=\gamma'/\|\gamma'\|$, $F_{\parallel}=F\cdot\tau$ is scalar.\\[3pt]
Only the tangential component contributes:\\[3pt]
$F\cdot\gamma'\Delta t=F_{\parallel}\|\gamma'\|\Delta t$.\\[5pt]
Length uses $\|\gamma'\|\Delta t$,\\
independent of the force or its sign.};

\end{tikzpicture}
\caption{The dot product extracts the force component along displacement. The blue vector follows the velocity $\gamma'(t)$ and approximates the actual short displacement along the curve.}
\label{fig:work-projection}
\end{figure}
